\documentclass[nonacm,manuscript,screen,microtype={protrusion=true,expansion,kerning,spacing}]{acmart}
\usepackage{amsmath}
\usepackage{mathtools}
\usepackage[english]{babel}
\usepackage{graphicx, lettrine}
\usepackage{calc}
\usepackage{bm}
\usepackage{epigraph}

\usepackage{pgfplots}
\pgfplotsset{compat=1.18}
\usetikzlibrary{backgrounds}

\begin{document}

\title{Editorial: Past, Present, and Future}
\subtitle{``EachMovie: archaeology of a first latent recommender system''}
\author{John DeTreville}
\email{john@detreville.org}

\maketitle{}

\section*{Introduction}

My ACM RecSys 2026 paper ``EachMovie: archaeology of a first latent recommender system''~\cite{DeTreville2026EachMovie_paper} is part of that conference's decennial ``Past, Present \& Future'' track. My conference presentation includes extra bonus material not in the paper.\footnote{The final presentation video of 2026-09-30 is expected to take some time to replace the ACM Digital Library's backup video from 2026-07-31.} This Editorial briefly lists some of the new points introduced there.

\emph{Warning}: The conference presentation includes this new material in a deliberately personal Editorial context. That same personal perspective continues here. This is all just my opinion. More research is clearly needed.

\section{Past}

As noted in the paper, EachMovie cultivated a high degree of user engagement and trust. Its users \emph{wanted} to use EachMovie, and to help it model their tastes and their souls.\footnote{The reference to ``souls'' here is deliberate. Movies, done well, reflect the human soul, and they shape the human soul. This is true of all good art. I posit that latent recommender systems for humans have a similar \emph{potential} to reflect the human soul (via their users' votes as carried into its internal models) and to shape the human soul (via their recommendations to their users). I believe this is true regardless of our current systems' massive shortcomings.}

EachMovie's users were invited to write in, and many did~\cite{EachMovieMailbag}. One particular respondent wrote:

\begin{quote}
  \emph{\large I've been using it for two months, and it's eerily good at predicting what I'll like~\ldots{} One of my favorite games with EachMovie is to enter a bogus rating, just to see how it affects the predictions; when I revert to my true opinion, EachMovie recalculates all its predictions instantly.~\ldots{} A better predictor, and more fun to use.~\cite{vanLeunen1996_mailbag}}
\end{quote}

The manner in which this user twiddled EachMovie's knobs is like a textbook example of \emph{human agency}~\cite{Bandura2001Agency} in an online system~\cite{Sundar2020}. This human user (whom I recognized as a former colleague) could act intentionally; she could control her environment; and she could shape her experiences. In short, EachMovie \emph{let her drive.}

The EachMovie team did not plan for this most excellent outcome. \emph{Letting users drive} was an emergent property, but it was a very welcome one. This outcome was completely accidental, but I argue that this is exactly the type of outcome we should try to replicate in our modern recommender systems.

\subsection{Letting users drive}

\emph{Letting users drive} is the important contribution of this Editorial. But read on anyway.

\pagebreak

\section{Present}

That was then and this is now. Thirty years have passed. Those carefree halcyon days of the 1990s are long gone.

My historical paper's Conclusions section was called ``Back to the future'' and it was boring. It arguably followed the similarly anodyne conclusion of that movie~\cite{Zemeckis1985BTTF}. This new editorial content in the conference presentation was meant to be less boring, more closely following the plot of ``Back to the Future Part II''~\cite{Zemeckis1989BTTF2}.

In this light, our own current world can perhaps be viewed as a divergent timeline, as in the second movie. Biff~Tannen\footnote{Biff Tannen is the movie's bad guy, who gains excessive control over the dystopian ``Hell Valley'' by unscrupulous means.} is running the casino,\footnote{This is ``Biff Tannen's Pleasure Paradise Casino \& Hotel.''} and he's dialing the figurative \emph{doomscrollamine}\footnote{This very recent neologism is from Andrew Marlton~\cite{Marlton2026Doomscrollamine}. It combines ``doomscrolling'' with ``dopamine,'' the chemical that floods the brain following rewards, helping to create addiction via neuroplasticity. That is not a good combination.} all the way up to 11. Ours is arguably the wrong timeline, at least for recommender systems and their users. We're all doomed. Or maybe not. Let's hope not.

\section{Future}
\epigraph{\raggedleft{} Yes, there are two paths you can go by, \\ but in the long run / \\ There's still time to change the road you're on.}{\textit{Stairway to Heaven}, Robert Plant, 1971}

What to do? I am not sure, but it ``really makes me wonder''~\cite{LedZeppelin1971Stairway}. Perhaps the first step is to build our systems to help users drive again.

\section{Summary}

EachMovie let users drive. We should do that again.

\begin{acks}
My wife is the fabulous Laurie Shaw-DeTreville. Many thanks to Laurie, and to Christine Seawell (UCI~English) as well, for so quickly reviewing drafts of these appendices, and for so patiently listening to me tell them all about my work over and over so many decades.
\end{acks}

\bibliographystyle{ACM-Reference-Format}
\bibliography{Appendix1.bib}

\end{document}